\section{本章小结}
图像先有正确的数据类型、通道和尺寸，后续操作才有意义。读图后判空，转换色彩空间，再根据任务分割、滤波和形态学处理，最后用轮廓描述几何。应熟练保存中间图、检查像素与计数；特征匹配和具体 API 参数可查文档。ROS 图像转换还要保持编码、时间戳与坐标系。

\begin{problemset}
\item 基础：CMake 已找到 OpenCV，但头文件或链接仍失败，应检查哪些目标设置？
\item 基础：a 为 10×20 的 CV\_8UC3 图像，不考虑行填充时占多少字节？b=a 后修改 b 是否影响 a？
\item 基础：imread 返回空 Mat 时，能否直接调用 cvtColor？imwrite 怎样选择保存格式？
\item 基础：8 位 HSV 图像中 H=178 与 H=2 是否可能都是红色？灰度值恰为 127、阈值也为 127，使用普通二值阈值模式时输出多少？
\item 基础：孤立椒盐点优先试哪种滤波？开运算和闭运算的先后顺序分别是什么？
\item 基础：20×10 的实心像素矩形为什么 contourArea 常为 171 而非 200？
\item 基础：将 100×60 图像各边缩为 0.5，输出尺寸是多少？为什么缩放后绘图坐标也要调整？
\item 基础：ORB 描述子常用什么距离？匹配距离小是否足以确认两点属于同一目标？
\item 提高：收到 Image 消息后只有通道数，没有检查 encoding，直接按 BGR 使用有什么问题？
\item 实践：运行 Mat 深浅拷贝示例，把 clone 改成共享赋值，记录输出变化并解释。
\end{problemset}

\section{本章验收}
\begin{enumerate}
    \item 用 CMake 构建程序，读取检查者给出的图片或摄像头；展示读取失败时的提示及正常显示。
    \item 完成 HSV 分割、开闭运算和轮廓筛选，保存每个阶段的中间图；红色样本覆盖色环两端。
    \item 绘制旋转矩形并解释面积和长宽比阈值；现场修改 ROI，展示共享赋值与 clone 对处理结果的影响。
    \item 选做：把图像转换为 ROS2 消息并显示，检查编码和时间戳是否保持一致。
\end{enumerate}
